\documentclass[11pt]{article}
\usepackage{amsmath}
\usepackage{amssymb}
\usepackage{geometry}
\geometry{margin=1in}
\usepackage{booktabs}
\usepackage{enumitem}

\title{Convexity-Based Approximation and Scaled Reduced-Form Model \\
with Mass of Lower-Ranked Workers and Firm Wage Allocation}
\author{}
\date{}

\begin{document}

\maketitle

\section{Core Idea: Convexity Approximation in the Non-Stochastic Model}

In the original N-worker model, the lottery delivers a stochastic prize drawn from a Pareto distribution with heavy right tail. Low-position workers favor it because the utility function is convex below the reference point $r$ ($\gamma_\text{low} > 1$), implying
\[
\mathbb{E}[u(y)] > u(\mathbb{E}[y])
\]
for random $y$ with positive variance. Convexity amplifies rare large upside outcomes, creating effective risk-loving behavior despite a roughly mean-preserving lottery.

In the simplified non-stochastic model, we avoid repeated random draws. Instead, we replace stochastic lottery prizes with a deterministic equivalent that includes a linearized risk premium capturing the second-order welfare effect of variance via utility curvature.

The effective lottery benefit for a representative agent in class $w$ includes a mean pool return share plus a risk-adjustment term
\[
\rho_w = \frac{1}{2} u''(m_w) \cdot \widehat{\sigma}_w^2,
\]
where a simple proxy for virtual variance is
\[
\widehat{\sigma}_w^2 = c \cdot ((1 - \lambda_w) m_w)^2 \qquad (c > 0 \text{ calibration constant}).
\]

\section{Scaled Reduced-Form Model with Mass at the Bottom}

We consider a simplified two-class model with unequal population: one representative higher-ranked worker (mass 1) with income $m_H = R (1-\mu)$, and $W \gg 1$ identical lower-ranked workers, each with per-capita income
\[
m_L = \frac{R \mu}{W}.
\]
Each representative agent in class $w$ chooses $\lambda_w \in [0,1]$, the **fraction of per-capita income allocated to upskilling** (safe improvement channel). The fraction allocated to the shared lottery pool is therefore $1 - \lambda_w$.

The upskilling cost for a representative agent is linear in the upskilling share:
\[
C_L(\lambda_L) = b \eta \lambda_L m_L, \qquad C_H(\lambda_H) = b \eta \lambda_H m_H,
\]
where $b > 1$ approximates the original power-law friction, and $\eta \in (0,1)$ is upskilling efficiency.

Aggregate lottery pool:
\[
\Pi = W (1 - \lambda_L) m_L + (1 - \lambda_H) m_H.
\]
Lottery return to a representative lower-ranked worker:
\[
\kappa W (1 - \lambda_L) m_L.
\]
Lottery return to the higher-ranked worker:
\[
\kappa (1 - \lambda_H) m_H.
\]

Net income per representative lower-ranked worker:
\begin{align*}
y_L &= m_L + \eta \lambda_L m_L - C_L(\lambda_L) - (1 - \lambda_L) m_L \\
    &\quad + \kappa W (1 - \lambda_L) m_L + \rho_L(1 - \lambda_L),
\end{align*}
with $\rho_L(1 - \lambda_L) = p_L (1 - \lambda_L)^2$ ($p_L > 0$).

Net income for the higher-ranked worker:
\begin{align*}
y_H &= m_H + \eta \lambda_H m_H - C_H(\lambda_H) - (1 - \lambda_H) m_H \\
    &\quad + \kappa (1 - \lambda_H) m_H + \rho_H(1 - \lambda_H),
\end{align*}
with $\rho_H(1 - \lambda_H) = p_H (1 - \lambda_H)^2$ ($p_H < 0$).

With linear utility $u(y) = y$, each maximizes own net income (Nash, treating other's choice fixed).

\subsection{Low-class first-order condition (per representative low worker)}
\[
\frac{\partial y_L}{\partial \lambda_L} = (\eta - b\eta - \kappa W + 1) m_L + 2 p_L (1 - \lambda_L) = 0
\]
Solving:
\[
\lambda_L^* = 1 - \frac{m_L (\kappa W - \eta + b\eta - 1)}{2 p_L},
\]
bounded as $\lambda_L^* = \max\!\bigl(0, \min\!\bigl(1, \text{above}\bigr)\bigr)$.

\subsection{High-class first-order condition}
\[
\frac{\partial y_H}{\partial \lambda_H} = (\eta - b\eta - \kappa + 1) m_H + 2 p_H (1 - \lambda_H) = 0
\]
Solving:
\[
\lambda_H^* = 1 - \frac{m_H (\kappa - \eta + b\eta - 1)}{2 p_H},
\]
bounded as $\lambda_H^* = \max\!\bigl(0, \min\!\bigl(1, \text{above}\bigr)\bigr)$.

\section{The Firm Layer: Wage Allocation and Asset Optimization}

The firm controls the total wage bill $R_t = \pi A_t$ (portion of current assets paid as wages). The firm chooses $\alpha \in [0,1]$ such that:
- Aggregate wages to lower-ranked workers (mass $W$) = $\alpha R_t$,
- Wages to the higher-ranked worker (mass 1) = $(1-\alpha) R_t$.

Per-capita wages are then $\alpha R_t / W$ for each lower-rank worker and $(1-\alpha) R_t$ for the higher-rank worker.

The firm maximizes the \emph{effective upskilling fraction} of the wage bill:
\[
\phi(\alpha) = \alpha \lambda_L^*(\alpha) + (1-\alpha) \lambda_H^*(\alpha),
\]
where $\lambda_L^*(\alpha)$ and $\lambda_H^*(\alpha)$ are the workers' equilibrium upskilling shares given policy $\alpha$.

Asset dynamics:
\begin{align*}
A_{t+1} &= (1 - \pi - \delta) A_t + \pi A_t \cdot \phi(\alpha) \\
        &\quad - C(A_t) - K(\alpha_t - \alpha_{t-1}),
\end{align*}
with $C(A_t)$ (maintenance cost, e.g. $c A_t^\gamma$, $\gamma > 1$) and $K(\Delta\alpha)$ (adjustment cost of changing $\alpha$, e.g. $k |\Delta\alpha|^\beta$, $\beta > 1$).

The firm chooses $\alpha_t$ to maximize $A_{t+1}$ (or long-run asset value), anticipating workers' responses. This creates a two-level Nash equilibrium:
- Workers optimize $\lambda_L$, $\lambda_H$ given $\alpha$.
- Firm optimizes $\alpha$ given $\lambda_L^*$, $\lambda_H^*$.

The effective fraction $\phi(\alpha)$ is the share of wages that becomes productive investment. The firm typically sets high $\alpha^*$ (large wage share to lower ranks) when $W$ is large, because the mass amplifies even modest upskilling from the bottom into large aggregate growth.

\subsection{Steady-State Equilibrium Analysis}

A steady-state equilibrium is a fixed point \((\alpha^*, \lambda_L^*, \lambda_H^*, A^*)\) satisfying two conditions simultaneously:

\begin{enumerate}
\item Given \(\alpha^*\), the workers choose their best-response upskilling shares:
\[
\lambda_L^* = 1 - \frac{m_L (\kappa W - \eta + b\eta - 1)}{2 p_L}, \qquad
\lambda_H^* = 1 - \frac{m_H (\kappa - \eta + b\eta - 1)}{2 p_H}
\]
(bounded to \([0,1]\)).
\item Given \(\lambda_L^*\) and \(\lambda_H^*\), the firm chooses \(\alpha^*\) to maximise the effective upskilling fraction
\[
\phi(\alpha) = \alpha \lambda_L^*(\alpha) + (1-\alpha) \lambda_H^*(\alpha).
\]
The steady-state asset level \(A^*\) then satisfies
\[
A^* = \frac{\pi \phi(\alpha^*)}{\delta + c (A^*)^{\gamma-1}},
\]
where \(C(A) = c A^\gamma\) (\(\gamma > 1\)).

The equilibrium is solved numerically by fixed-point iteration: guess \(\alpha_0\), compute \(\lambda_L^*\) and \(\lambda_H^*\), update \(\alpha\) to maximise \(\phi(\alpha)\), and repeat until convergence. Convergence is fast and stable.

\subsubsection{Comparative Statics}

The table below summarises directional effects of key parameters on the equilibrium values (holding others fixed).

\begin{table}[ht]
\centering
\begin{tabular}{lccc}
\toprule
Parameter & \(\alpha^*\) (firm) & \(\lambda_L^*\) (low-rank) & \(\lambda_H^*\) (high-rank) \\
\midrule
\(\uparrow W\) (larger mass of poor) & \(\uparrow\) & \(\downarrow\) & \(\approx 1\) (corner) \\
\(\uparrow \eta\) (better upskilling) & \(\uparrow\) or \(\downarrow\) (ambiguous) & \(\uparrow\) & \(\uparrow\) \\
\(\uparrow p_L\) (stronger low-rank lottery premium) & \(\downarrow\) & \(\downarrow\) & little change \\
\(\uparrow |p_H|\) (stronger high-rank aversion) & \(\uparrow\) & little change & \(\uparrow\) \\
\(\uparrow \pi\) (higher wage share) & \(\uparrow\) & little direct effect & little direct effect \\
\(\uparrow \delta\) (higher depreciation) & little direct effect & \(\uparrow\) (indirect) & \(\uparrow\) (indirect) \\
\bottomrule
\end{tabular}
\caption{Directional comparative statics at equilibrium.}
\end{table}

\subsubsection{Equilibrium Implications}

The model produces several robust implications:

\begin{itemize}
\item \textbf{Mass effect dominates}: Large \(W\) pushes the firm toward high \(\alpha^*\) (allocating a large wage share to lower ranks) to leverage the numerous workers. Even modest improvements in \(\lambda_L\) generate large aggregate upskilling.
\item \textbf{Sticky top, churning bottom}: The high-rank worker typically stays at \(\lambda_H^* \approx 1\) (full upskilling), reinforcing stability at the top. Lower-rank workers allocate substantially to the lottery (\(\lambda_L^* < 1\)), producing persistent churning at the bottom.
\item \textbf{Endogenous growth channel}: Unlike the classic Solow model (exogenous savings rate \(s\)), here the effective investment rate \(\pi \phi(\alpha^*)\) is chosen strategically by the firm. Growth \(g^* = \pi \phi(\alpha^*) - \delta - C(A^*)/A^*\) therefore depends on policy responsiveness and population structure.
\item \textbf{Policy relevance}: The firm has a meaningful lever (\(\alpha\)) to counteract baseline inequality (\(\mu\)). When \(W\) is large, optimal policy often favours the bottom, illustrating how institutional wage allocation can influence long-run inequality and growth.
\item \textbf{Two-level Nash character}: The equilibrium is a mutual best-response fixed point. Small changes in parameters (e.g., stronger risk premium \(p_L\)) can shift the entire system through feedback between workers and firm.
\end{itemize}

Overall, the model shows that when lower ranks are numerous (\(W \gg 1\)), rational firm behaviour naturally tilts toward high \(\alpha^*\), boosting aggregate growth while preserving the micro-level asymmetry of churning at the bottom and stickiness at the top.
\item 
\section{New Model}


\end{document}
\end{document}